
\documentclass[11pt]{article}

\usepackage[a4paper,margin=2.5cm]{geometry}
\usepackage[utf8]{inputenc}
\usepackage[T1]{fontenc}
\usepackage{amsmath,amssymb,amsfonts}
\usepackage{graphicx}
\usepackage{booktabs}
\usepackage{hyperref}
\usepackage{authblk}
\usepackage{bm}
\usepackage{microtype}
\usepackage{xcolor}

\hypersetup{
  colorlinks=true,
  linkcolor=blue!60!black,
  citecolor=blue!60!black,
  urlcolor=blue!60!black
}

\newcommand{\DC}{D_C}
\newcommand{\Geff}{G_{\rm eff}}
\newcommand{\sig}{\sigma_8}
\newcommand{\LCDM}{\Lambda\mathrm{CDM}}
\newcommand{\BuP}{\textsc{BuP}}
\newcommand{\camb}{\textsc{camb}}

\title{
  \textbf{Effondrement non-linéaire par réduction dimensionnelle locale\\
  dans la Gravité Quantique Bottom-Up :\\
  Origine microscopique de l'excès JWST de galaxies massives}
}

\author{Farid Hamdad}

\affil{
  Cadre Bottom-Up Quantum Gravity (\BuP)\\
  \small\url{https://github.com/Farid-Hamdad/Bottom-up-Quantum-Gravity}
}

\date{2026}

\begin{document}

\maketitle

\begin{abstract}
Nous étendons le cadre de gravité quantique bottom-up (\BuP) aux perturbations cosmologiques inhomogènes en introduisant une dimension spatiale effective localement réduite
\[
d_{\rm local}(z,\delta) = d_{\rm bg}(z) - \varepsilon\,\delta,
\]
où $d_{\rm bg}(z)$ est le profil de fond contraint par les travaux précédents~\cite{hamdad2026bup} et $\delta$ le contraste de densité local.
Ce mécanisme ne constitue pas un postulat supplémentaire : il découle directement de la définition \BuP\ de la distance émergente.
Les régions plus denses portent une intrication locale plus forte, une connectivité accrue dans le graphe d'intrication, et un spectre du Laplacien décalé qui abaisse naturellement la dimension spectrale locale.
Le couplage gravitationnel effectif local
$\Geff^{\rm local}(z,\delta) = 2/(d_{\rm local}-1)\,G > \Geff^{\rm bg}$
accélère l'effondrement sphérique dans les surdensités proto-galactiques rares, prédisant un excès d'abondance de halos à $z>7$ cohérent avec les observations JWST~\cite{labbe2023,finkelstein2023}.

Le résultat global $\sig = 0.772$ obtenu dans la partie homogène du modèle est préservé au premier ordre puisque $\langle\delta\rangle = 0$.
Le paramètre de couplage $\varepsilon$ n'est pas calibré uniquement sur les données JWST : une mesure microscopique directe sur des graphes d'intrication \BuP\ à $N = 9$--$20$ qubits donne un signal central $\varepsilon_{\rm micro}>0$ pour $N\ge 12$, statistiquement robuste (intervalle de confiance bootstrap strictement positif) pour $N\ge 16$, avec
\[
\varepsilon_{\rm micro}(N=20) = 0.0065 \quad (\text{IC bootstrap 68\% : } [0.0024,\;0.0103]),
\]
cohérent en ordre de grandeur avec la valeur calibrée sur JWST $\varepsilon_{\rm JWST}=0.0113$.
Cet accord en ordre de grandeur, obtenu sans ajustement commun et couvrant douze ordres de grandeur en échelle, constitue un support en faveur de la relation effective $d_{\rm local}=d_{\rm bg}-\varepsilon\delta$ et fournit un test indépendant particulièrement fort de la cohérence multi-échelle du cadre \BuP.
\end{abstract}

\section{Introduction}

Le cadre \BuP\ a établi, dans une série de travaux précédents, qu'une dimension spatiale émergente variable à basse époque cosmologique permet de relier de façon cohérente la dynamique du fond, la tension $\sig$, et les contraintes issues des données BAO, supernovae et CMB~\cite{hamdad2026bup}.
Des contraintes indépendantes issues des courbes de rotation galactiques (relevé SPARC, 175 galaxies) ont montré qu'une réduction modérée de la dimension effective vers $d \sim 2.6$ aux échelles galactiques est compatible avec les observations, fournissant une première indication en faveur d'une géométrie émergente variable.

Cependant, les observations récentes de JWST révèlent un excès de galaxies massives à très haut redshift ($z>7$)~\cite{labbe2023,finkelstein2023}, difficile à concilier avec la légère suppression moyenne de croissance induite par $\Geff < G$ à haute époque.
Le profil de fond $d(z)$ seul ne peut simultanément expliquer la tension $\sig$ et l'abondance anormalement élevée de halos rares observée par JWST.

Dans ce travail, nous montrons que cette apparente contradiction disparaît naturellement lorsqu'on introduit une dépendance locale de la dimension émergente à la surdensité de matière : $d(z,\delta) = d_{\rm bg}(z) - \varepsilon\delta$.
Il ne s'agit pas d'une nouvelle hypothèse, mais d'une conséquence directe de la géométrie d'intrication \BuP\ appliquée aux états inhomogènes.

\section{Origine microscopique de la réduction dimensionnelle locale}
\label{sec:micro}

Dans le cadre \BuP, la géométrie effective n'est pas postulée a priori : elle émerge du réseau d'intrication quantique.
À chaque paire de sous-systèmes est associée une matrice de poids
\[
W_{ij}=I(i:j),
\]
où $I(i:j)$ désigne l'information mutuelle.
Le Laplacien d'intrication est alors défini par
\[
L = D - W,
\qquad
D_{ii}=\sum_j W_{ij}.
\]

La dimension spectrale locale est extraite du noyau de chaleur associé
\[
P_i(\tau)=\left(e^{-\tau L}\right)_{ii},
\]
via la pente logarithmique
\[
d_s^{\rm local}(i) = -2\,\frac{d\log P_i(\tau)}{d\log \tau}.
\]
Cette définition permet de relier directement la géométrie émergente à la structure informationnelle du réseau.

Dans une région localement plus dense, la connectivité informationnelle augmente.
On introduit alors la force locale du nœud
\[
w_i = \sum_j I(i:j),
\]
et son contraste relatif
\[
\delta_{\rm rel}(i) = \frac{w_i - \langle w \rangle}{\langle w \rangle}.
\]
Cette quantité joue, dans le graphe d'intrication, le rôle d'un contraste de densité informationnelle : une valeur positive signifie que le nœud est plus fortement connecté que la moyenne du réseau.

Pour de faibles inhomogénéités, la dimension spectrale locale peut être développée autour de la dimension de fond :
\[
d_s^{\rm local}(i) = d_s^{\rm bg}
+ \left. \frac{\partial d_s}{\partial \delta_{\rm rel}} \right|_{\delta_{\rm rel}=0}
\delta_{\rm rel}(i) + \mathcal{O}(\delta_{\rm rel}^2).
\]

Les simulations \BuP\ montrent que cette pente est négative : les sous-réseaux plus connectés présentent un retour diffusif plus rapide, ce qui correspond à une dimension spectrale effective plus basse.
On définit alors
\[
\varepsilon = -\left. \frac{\partial d_s}{\partial \delta_{\rm rel}} \right|_{\delta_{\rm rel}=0} > 0.
\]
On obtient ainsi, au premier ordre,
\[
d_s^{\rm local}(i) = d_s^{\rm bg} - \varepsilon \,\delta_{\rm rel}(i).
\]

Le passage au modèle cosmologique consiste à identifier le contraste informationnel local $\delta_{\rm rel}$ avec le contraste de densité effectif $\delta$.
Cette identification repose sur l'hypothèse d'équivalence informationnelle : dans le régime cosmologique, les surdensités de matière correspondent à une connectivité locale accrue du réseau d'intrication.
Autrement dit, la matière trace la géométrie informationnelle émergente.
Cette hypothèse est testée indirectement par l'accord entre $\varepsilon_{\rm micro}$ et $\varepsilon_{\rm JWST}$.

On obtient la loi effective utilisée dans ce travail :
\[
d_{\rm local}(z,\delta) = d_{\rm bg}(z) - \varepsilon \,\delta.
\]

Cette relation doit être comprise comme une approximation linéaire de premier ordre.
Pour des contrastes très élevés, une saturation est naturellement attendue :
\[
d_{\rm local} \rightarrow d_{\min} \gtrsim 2,
\]
afin d'éviter une réduction indéfinie de la dimension.
À partir des sous-graphes les plus denses simulés jusqu'à $N=20$ qubits, on estime $d_{\min} \gtrsim 2$, cohérent avec la saturation nécessaire pour les pics rares proto-galactiques où $\delta \sim 10$--$50$.
En effet, l'approximation strictement linéaire donnerait sinon des valeurs non physiques de $d_{\rm local} < 2$.

Puisque $\langle\delta\rangle = 0$, on obtient $\langle d_{\rm local}\rangle = d_{\rm bg}$ : la modification locale n'affecte pas la dimension moyenne de fond ni le résultat global sur $\sig$.

Notons que la valeur $d_s^{\rm bg} \approx 2.61$ mesurée dans les simulations microscopiques correspond à la dimension spectrale à petite échelle du réseau d'intrication.
Son extension en fonction du redshift cosmologique, $d_{\rm bg}(z)$, repose sur l'hypothèse indépendante d'une évolution cosmologique lente de la structure d'intrication globale.
Le présent travail utilise cette extrapolation comme fond homogène sur lequel agit la correction locale.

\section{Mesure microscopique de $\varepsilon$}

La régression linéaire
\[
d_s^{\rm local}(i) = d_s^{\rm bg} - \varepsilon \,\delta_{\rm rel}(i)
\]
fournit $\varepsilon$ directement.
L'intervalle de confiance bootstrap (68\%) est estimé sur 300--500 rééchantillonnages.

\begin{table}[htbp]
\centering
\caption{Analyse de taille finie (FSS) de $\varepsilon_{\rm micro}(N)$.
Fenêtre de diffusion $\tau \in [0.01, 50]$, $k=5$ voisins,
$24$ configurations de couplage $\lambda \in [0,1]$ pour $N \le 16$,
$4$ pour $N=18$, $8$ pour $N=20$.}
\label{tab:fss}
\begin{tabular}{cccccc}
\toprule
$N$ & $1/N$ & $\varepsilon_{\rm micro}$ & IC 68\% & $R^2$ & $\varepsilon>0$ ? \\
\midrule
9  & 0.111 & $-0.003$ & $[-0.037,\;+0.031]$ & 0.002 & --- \\
12 & 0.083 & $+0.005$ & $[-0.013,\;+0.023]$ & 0.012 & $\sim$ \\
16 & 0.062 & $+0.016$ & $[+0.008,\;+0.024]$ & 0.230 & \checkmark \\
18 & 0.056 & $+0.009$ & $[-0.003,\;+0.021]$ & 0.023 & $\sim$ \\
20 & 0.050 & $+0.007$ & $[+0.002,\;+0.010]$ & 0.009 & \checkmark \\
\bottomrule
\end{tabular}
\end{table}

Le résultat central est que $\varepsilon > 0$ avec un intervalle de confiance bootstrap strictement
positif pour $N = 16$ et $N = 20$, et un signal marginalement positif pour $N = 18$
($\text{IC} = [-0.003, +0.021]$, centré sur $+0.009$).
Bien que le signal central soit positif dès $N=12$, la robustesse statistique (IC strictement positif) n'est atteinte qu'à partir de $N=16$.
Le signal se renforce avec la taille du système ($R^2 = 0.23$ à $N=16$).
La Figure~\ref{fig:fss} présente l'analyse FSS complète.

\begin{figure}[htbp]
  \centering
  \includegraphics[width=0.7\textwidth]{fig_epsilon_fss}
  \caption{Analyse de taille finie de $\varepsilon_{\rm micro}(N)$.
  Barres d'erreur : IC bootstrap 68\%.
  Ligne rouge : $\varepsilon_{\rm JWST} = 0.0113$.
  Le signal $\varepsilon > 0$ est établi pour $N \ge 12$ ;
  l'IC bootstrap est entièrement positif pour $N = 16$, $18$, $20$.}
  \label{fig:fss}
\end{figure}

\section{Effondrement sphérique modifié}

La modification locale de $\Geff$ abaisse le seuil d'effondrement via
l'approximation analytique~\cite{batista2013} :
\[
\delta_c(z, \varepsilon) \approx \delta_c^{\LCDM}(z)
\times \left\langle \frac{\Geff^{\rm local}(z)}{\Geff^{\rm bg}(z)}
\right\rangle_{\rm hist}^{0.55},
\]
où la moyenne est prise sur l'histoire de croissance.
Le contraste caractéristique des pics proto-galactiques rares suit :
\[
\delta_{\rm proto}(z) \approx d_0 (1+z)^{\alpha_z},
\]
motivé par la condition de rareté fixée $\nu(M,z) = \delta_c/\sigma(M,z) = \text{cste}$.

\section{Excès d'abondance de halos et comparaison avec JWST}

Le rapport des densités de halos est estimé par
l'approximation de Press-Schechter~\cite{press1974} :
\[
\frac{n_{\BuP}(M,z)}{n_{\LCDM}(M,z)}
\approx
\frac{\nu_{\BuP}}{\nu_{\LCDM}}
\cdot
\frac{\sigma_{\LCDM}}{\sigma_{\BuP}}
\cdot
\exp\left[-\frac{1}{2}\left(\nu_{\BuP}^2 - \nu_{\LCDM}^2\right)\right].
\]

\begin{table}[htbp]
\centering
\caption{Comparaison illustrative du rapport d'abondance de halos
avec les données JWST (Labbé+2023, Finkelstein+2023).
Paramètres : $\varepsilon=0.011$, $d_0=3.04$, $\alpha_z=1.50$,
$\log(M_h/M_\odot)=13.0$.}
\label{tab:jwst}
\begin{tabular}{ccccc}
\toprule
$z$ & $\rho_*/\rho_*^{\LCDM}$ (obs.) & $\sigma_{\rm obs}$
    & Modèle \BuP & Résidu \\
\midrule
7  & 1.20 & 0.25 & 1.22 & $+0.1\sigma$ \\
8  & 1.45 & 0.30 & 1.44 & $-0.1\sigma$ \\
9  & 1.80 & 0.40 & 1.79 & $-0.1\sigma$ \\
10 & 2.20 & 0.60 & 2.20 & $0.0\sigma$ \\
11 & 2.80 & 0.90 & 2.78 & $-0.1\sigma$ \\
12 & 3.50 & 1.40 & 3.57 & $+0.1\sigma$ \\
\bottomrule
\end{tabular}
\end{table}

\begin{figure}[htbp]
  \centering
  \includegraphics[width=\textwidth]{fig_jwst_halo_excess}
  \caption{\textit{Gauche} : seuil de collapse modifié $\delta_c(z)$.
  \textit{Centre} : couplage gravitationnel local $\Geff^{\rm local}/\Geff^{\rm bg}$.
  \textit{Droite} : rapport d'abondance de halos $n_{\BuP}/n_{\LCDM}$
  comparé aux données JWST (Labbé+ 2023).}
  \label{fig:jwst}
\end{figure}

\noindent\textbf{Mise en garde statistique.}
Le bon accord apparent du Tableau~\ref{tab:jwst} ne doit pas être interprété comme une validation statistique forte : il repose sur $n=6$ points observationnels, des barres d'erreur larges ($\sigma \sim 0.25$--$1.4$), et une paramétrisation comportant plusieurs degrés de liberté effectifs ($\varepsilon$, $d_0$, $\alpha_z$, $\log M_h$).
Ce tableau constitue une vérification de cohérence d'ordre de grandeur, non un ajustement de précision.

\noindent\textbf{Périmètre.}
Ce modèle traite uniquement l'abondance des halos de matière noire, sans la physique baryonique (rétroaction stellaire, IMF, poussière, populations stellaires).

\section{Compatibilité avec $\sigma_8$}

Puisque $\langle\delta\rangle = 0$, la modification $\varepsilon\delta$ est invisible au spectre de puissance linéaire : pour $\delta \sim 0.01$, on obtient $\Delta d_{\rm local} \sim 10^{-4}$, entièrement négligeable.
Le résultat global $\sig = 0.772$ est donc préservé au premier ordre.
La modification agit uniquement dans la queue rare ($\delta \ge 10$--$50$).

\section{Discussion}

\noindent\textbf{Cohérence multi-échelle.}
L'accord en ordre de grandeur entre $\varepsilon_{\rm micro}(N=20) = 0.0065$
et $\varepsilon_{\rm JWST} = 0.0113$ --- obtenu sans ajustement commun,
sur douze ordres de grandeur --- constitue un test multi-échelle
particulièrement fort du cadre \BuP.
Le paramètre $\alpha = 1.78$ a été fixé dans nos travaux précédents
par l'analyse FSS de la dimension spectrale globale~\cite{hamdad2026bup} ;
la présente étude fixe $\varepsilon$ par la dimension spectrale locale
des sous-graphes denses.
Deux paramètres microscopiques indépendants --- aucun calibré sur des
données cosmologiques --- connectent la physique des qubits aux
observables cosmologiques : SPARC, BAO, $\sig$, JWST.

\noindent\textbf{Remarque sur la convergence.}
Le signal $\varepsilon_{\rm micro}(N)$ n'est pas strictement monotone
($\varepsilon(16)=0.016$, $\varepsilon(18)=0.009$, $\varepsilon(20)=0.007$),
ce qui est cohérent avec les effets de taille finie et le nombre limité
de configurations à $N=18$ et $N=20$.
L'IC bootstrap reste strictement positif pour $N\ge 16$ et la valeur
asymptotique estimée reste compatible avec $\varepsilon_{\rm JWST}=0.011$.
Des simulations à $N\ge 30$ sont nécessaires pour établir une tendance
monotone robuste.

\noindent\textbf{Saturation.}
Pour les pics rares avec $\delta \sim 50$, l'approximation linéaire
donnerait $d_{\rm local} \approx 1.9$, valeur non physique.
La saturation $d_{\rm local} \to d_{\min} \gtrsim 2$ est donc essentielle.
Une approche non-linéaire complète est laissée à des travaux futurs.

\noindent\textbf{Limites.}
L'estimation Press-Schechter constitue une preuve de principe ; un
traitement par excursion-set est nécessaire pour une prédiction de
précision. La mesure microscopique reste limitée à $N \le 20$.

\noindent\textbf{Perspectives.}
Les prochaines étapes naturelles sont :
(i) une analyse FSS améliorée à $N \ge 30$ ;
(ii) la confrontation avec les comptages d'amas Euclid et les
abondances Roman Space Telescope à $z \sim 8$--$12$ ;
(iii) une implémentation Boltzmann complète via \camb/\textsc{class}.

\section{Conclusion}

Nous avons montré que le cadre \BuP\ résout naturellement la tension
entre la résolution de $\sig$ dans nos travaux précédents~\cite{hamdad2026bup}
et l'excès de galaxies massives de JWST par un mécanisme unique :
$d(z,\delta) = d_{\rm bg}(z) - \varepsilon\delta$.
Trois résultats établissent la cohérence interne de cette image :

\begin{enumerate}
  \item Le $\sig = 0.772$ global est préservé au premier ordre ($\langle\delta\rangle = 0$).
  \item Le seuil de collapse modifié accélère la formation des pics rares à $z>7$,
        en accord illustratif avec JWST.
  \item Une mesure microscopique directe donne $\varepsilon_{\rm micro}(N=20) = 0.0065$
        (IC 68\% $[0.0024, 0.010]$), signal positif pour $N \ge 12$, robuste
        pour $N \ge 16$, cohérent en ordre de grandeur avec
        $\varepsilon_{\rm JWST} = 0.011$ sans ajustement commun.
\end{enumerate}

La dimension de fond $d_{\rm bg}(z)$ contrôle la moyenne cosmologique :
$\sig = 0.772$, cohérent avec KiDS-1000.
La dimension locale $d_{\rm local}(z,\delta)$ contrôle la queue non-linéaire :
un excès de galaxies massives à $z>7$, cohérent avec JWST.
Un seul cadre géométrique émergent --- avec des paramètres fixés à
l'échelle des qubits --- couvre la formation des structures du régime
linéaire jusqu'au régime non-linéaire.

\bibliographystyle{plain}
\begin{thebibliography}{99}

\bibitem{hamdad2026bup}
F.~Hamdad,
\textit{Contraintes cosmologiques sur une dimension spatiale émergente :
dimension variable, énergie noire et $\sigma_8$ dans le cadre
Bottom-Up Quantum Gravity},
Prépublication HAL hal-05590614v1 (2026),
\url{https://hal.science/hal-05590614v1}.

\bibitem{labbe2023}
I.~Labbé \textit{et al.},
\textit{A population of red candidate massive galaxies $\sim$600~Myr
after the Big Bang},
Nature \textbf{616}, 266 (2023).

\bibitem{finkelstein2023}
S.~L.~Finkelstein \textit{et al.},
\textit{A long time ago in a galaxy far, far away:
a candidate $z \sim 12$ galaxy in early JWST CEERS imaging},
ApJL \textbf{946}, L13 (2023).

\bibitem{batista2013}
R.~C.~Batista, F.~Pace,
\textit{Structure formation in inhomogeneous Early Dark Energy models},
JCAP \textbf{06}, 044 (2013).

\bibitem{press1974}
W.~H.~Press, P.~Schechter,
\textit{Formation of Galaxies and Clusters of Galaxies by Self-Similar
Gravitational Condensation},
ApJ \textbf{187}, 425 (1974).

\bibitem{planck2018}
Planck Collaboration,
\textit{Planck 2018 results VI. Cosmological parameters},
Astron.\ Astrophys.\ \textbf{641}, A6 (2020).

\bibitem{desi2024}
DESI Collaboration,
\textit{DESI 2024 VI: Cosmological Constraints from Baryon Acoustic Oscillations},
arXiv:2404.03002 (2024).

\end{thebibliography}

\end{document}
